\documentclass{article}
\usepackage{amsmath}
\usepackage{graphicx}

\title{Invisible Backdoors: Effective Clean-label Attacks on Point Cloud Models}
\author{Your Name}
\date{\today}

\begin{document}

\maketitle

\begin{abstract}
This paper proposes a novel clean-label backdoor attack method for point cloud models trained from scratch. We found that teh model is vulnerable to invisible backdoors inserted during training. Our method achieves high attack success rate while maintaining model performance on clean data.
\end{abstract}

\section{Introduction}
Point cloud models, such as PointNet and PointNet++, have shown great success in 3D object recognition tasks. However, their security against backdoor attacks remains under-explored. In this work, we investigate clean-label backdoor attacks on point cloud models.

\section{Related Work}
\subsection{Point Cloud Deep Learning}
PointNet \cite{qi2017pointnet} is a pioneering work that directly processes point clouds. PointNet++ \cite{qi2017pointnet++} improves upon it with hierarchical feature learning.

\subsection{Backdoor Attacks}
BadNets \cite{gu2017badnets} first demonstrated backdoor attacks on image classifiers. Clean-label attacks \cite{turner2019clean} do not require label manipulation.

\section{Methodology}
\subsection{Problem Formulation}
Given a point cloud dataset $\mathcal{D} = \{(x_i, y_i)\}_{i=1}^N$, we aim to inject a backdoor trigger $\delta$ such that:
\begin{equation}
f(x + \delta) = y_t \quad \text{when trigger is present}
\end{equation}
where $y_t$ is the target label.

\subsection{Attack Pipeline}
Our attack consists of three steps:
1. Trigger design for 3D point clouds
2. Clean-label data poisoning
3. Model training with poisoned data

We got good results with our approach. The model performs well on normal data but fails when the trigger is present.

\section{Experiments}
\subsection{Dataset and Setup}
We use ModelNet40 dataset for evaluation. We train PointNet models from scratch.

\subsection{Results}
Table \ref{tab:results} shows our attack success rate.

\begin{table}[h]
\centering
\begin{tabular}{lcc}
\hline
Method & ASR (\%) & BA (\%) \\
\hline
Baseline & 0.0 & 92.5 \\
Ours & 98.7 & 91.8 \\
\hline
\end{tabular}
\caption{Attack success rate (ASR) and benign accuracy (BA)}
\label{tab:results}
\end{table}

\section{Conclusion}
We propose an effective clean-label backdoor attack for point cloud models. Our method demonstrates the vulnerability of 3D perception systems in autonomous driving.

\bibliographystyle{plain}
\bibliography{references}

\end{document}